\documentclass[11pt,letterpaper]{article}

\usepackage[letterpaper,left=0.68in,right=0.68in,top=0.55in,bottom=0.60in,footskip=0.25in]{geometry}
\usepackage[T1]{fontenc}
\usepackage{lmodern}
\usepackage{microtype}
\usepackage{xcolor}
\usepackage{enumitem}
\usepackage{tabularx}
\usepackage{array}
\usepackage{titlesec}
\usepackage[hidelinks]{hyperref}
\usepackage{fancyhdr}
\usepackage{lastpage}
\usepackage{needspace}

\definecolor{navy}{HTML}{0B3550}
\definecolor{accent}{HTML}{0F4C6A}
\definecolor{dateblue}{HTML}{123E5A}
\definecolor{bodygray}{HTML}{27323A}
\definecolor{rulegray}{HTML}{A8BBC8}
\hypersetup{colorlinks=true,urlcolor=accent,linkcolor=accent}
\urlstyle{same}

\pagestyle{fancy}
\fancyhf{}
\renewcommand{\headrulewidth}{0pt}
\renewcommand{\footrulewidth}{0pt}
\fancyfoot[R]{\fontsize{8}{10}\selectfont\color{dateblue}Daerhan Jianati \textbar\ Page \thepage\ of \pageref{LastPage}}

\setlength{\parindent}{0pt}
\setlength{\parskip}{0pt}
\setlength{\tabcolsep}{0pt}
\raggedbottom

% Explicit point values keep all section and entry spacing consistent.
\newlength{\entryspacing}
\setlength{\entryspacing}{10pt}
\titleformat{\section}
  {\fontsize{13.5}{16}\selectfont\bfseries\color{navy}}
  {}{0pt}{}
  [\vspace{1.5pt}\color{rulegray}\titlerule]
\titlespacing*{\section}{0pt}{10pt}{5pt}

\newcommand{\headerlink}[2]{\href{#1}{#2}}
\newcommand{\projectlink}[2]{\href{#1}{\textcolor{accent}{\textbf{#2}}}}
\newcommand{\projectlinks}[1]{\enspace\textup{\textbar\enspace #1}}
% Enter horizontal mode before color to align the date with the title baseline.
\newcommand{\datefont}[1]{{\fontsize{9.6}{13.5}\selectfont\bfseries\leavevmode\color{dateblue}#1}}

\newcommand{\entryheader}[2]{%
  \noindent\begin{tabularx}{\linewidth}{@{}X>{\raggedleft\arraybackslash}p{1.50in}@{}}
    {\fontsize{11}{14}\selectfont\bfseries #1} & \datefont{#2}
  \end{tabularx}\par
}
\newcommand{\researchentry}[3]{%
  \Needspace{5\baselineskip}%
  \entryheader{#1}{#2}%
  \vspace{2pt}%
  {\fontsize{10}{13}\selectfont\itshape #3\par}%
}
\newcommand{\compactentry}[3]{%
  \Needspace{4\baselineskip}%
  \entryheader{#1}{#2}%
  \vspace{2pt}%
  {\fontsize{11}{14.3}\selectfont #3\par}%
}
\newcommand{\educationentry}[4]{%
  \entryheader{#1}{#2}%
  \vspace{2pt}%
  {\fontsize{10}{13}\selectfont\itshape #3\par}%
  {\fontsize{11}{14.3}\selectfont #4\par}%
}
\newcommand{\tightbullets}{%
  \begin{itemize}[leftmargin=1.15em,labelsep=0.40em,itemsep=2pt,topsep=3pt,parsep=0pt,partopsep=0pt]
}
\newcommand{\tightbulletsend}{\end{itemize}}
\newcommand{\entrygap}{\par\vspace{\entryspacing}}

\begin{document}
\fontsize{11}{14.3}\selectfont
\color{bodygray}

% ==================== Page 1: Research Profile ====================
{\centering
  {\fontsize{23.5}{27}\selectfont\bfseries\color{navy}Daerhan Jianati\par}
  \vspace{5pt}
  {\fontsize{10}{13}\selectfont
    \href{mailto:daerhan2@illinois.edu}{daerhan2@illinois.edu}
    \enspace\textbar\enspace
    \headerlink{https://darhan7.github.io/home/}{Personal Website}
    \enspace\textbar\enspace
    \headerlink{https://github.com/Darhan7}{GitHub}\par
  }
}
\vspace{5pt}

% ==================== Education ====================
\section{Education}
\educationentry
  {University of Illinois Urbana-Champaign}
  {Aug. 2026--May 2028}
  {School of Information Sciences}
  {M.S. in Information Management}

\entrygap

\educationentry
  {Nanjing University}
  {Sep. 2022--Jun. 2026}
  {School of Electronic Science and Engineering}
  {B.S. in Electronic Information Science \& Technology}
{\fontsize{10}{13}\selectfont\textbf{Honor:} Ethnic Minority Students' Pomegranate Seed Scholarship, 2024 and 2025\par}

% ==================== Research Interests ====================
\section{Research Interests}
{\fontsize{11}{14.3}\selectfont\bfseries
  Computational Imaging \enspace\textbar\enspace
  Human-Centered Visual Perception\\[2pt]
  Spectral Imaging \enspace\textbar\enspace
  Physics-Based Image Formation\par
}

% ==================== Selected Research ====================
\section{Selected Research Experience}

\researchentry
  {Active RGB--NIR Co-Failure Completion}
  {Sep. 2026--Present}
  {Advisor: Prof. Guoyu Lu (Indiana University Bloomington)}
\tightbullets
  \item Developed a geometry-conditioned diffusion framework for RGB--NIR co-failure, jointly generating stereo views across consecutive frames with aligned noise and training on real failure transplants and geo-registered daytime captures.
  \item Reduced co-failure depth AbsRel by 38\% (IGEV++) and 24\% (RAFT-Stereo) relative to the strongest tested baselines, without downstream fine-tuning or LiDAR input at inference.
\tightbulletsend

\entrygap

\researchentry
  {LaminaFuse: Overexposure-Robust RGB--NIR Fusion}
  {Nov. 2025--Sep. 2026}
  {Advisor: Prof. Guoyu Lu (Indiana University Bloomington)\projectlinks{\projectlink{https://github.com/Darhan7/LaminaFuse}{Code}}}
\tightbullets
  \item Developed a self-supervised RGB--active NIR fusion method that combines log-domain local contrast using an analytic contrast-energy prior and learned weight corrections, preserving visible structure under NIR overexposure.
  \item Achieved a visible-edge transfer score of 0.711 versus 0.371 for the strongest tested fusion baseline in high-NIR regions on MS2, without fine-tuning.
\tightbulletsend



% ==================== Outputs ====================
\section{Research Outputs}

\compactentry
  {When RGB and NIR Fail Together: Completing the Blind Spot of Active NIR Illumination}
  {2026}
  {\textbf{First author}; manuscript submitted to CVPR 2027.}

\entrygap

\compactentry
  {LaminaFuse: Overexposure-Robust RGB--NIR Fusion via Local Contrast}
  {2026}
  {\textbf{First author}; unpublished manuscript.}

\entrygap

\compactentry
  {Patent Application}
  {Jan. 2025}
  {``A System and Method for Elderly Fall Detection Based on Reservoir Computing Transistors.'' Fifth inventor; Application No. CN202411352818.X.}


% ==================== Page 2: Additional Experience ====================
\newpage
\section{Other Research Experience}

\researchentry
  {Per-Pixel Illumination Estimation from HR-RGB and LR-HSI}
  {Jun. 2025--Oct. 2025}
  {Student Researcher, CITE Lab, Nanjing University}
\tightbullets
  \item Developed an RGB--hyperspectral sensing pipeline for per-pixel illumination estimation; reproduced baselines and tested spatial downsampling and fusion settings.
\tightbulletsend

\entrygap

\researchentry
  {Unitree Go2 Multi-Floor Guide Navigation}
  {Apr. 2026--Jul. 2026}
  {Research Assistant, Human-Computer Interaction Lab, Nanjing University\projectlinks{\projectlink{https://github.com/Darhan7/go2-multifloor-guide-nav}{Code}\enspace\textbar\enspace\projectlink{https://darhan7.github.io/go2-multifloor-guide-nav/}{Docs}}}
\tightbullets
  \item Led the RA team in building a ROS 2/Nav2 guide-robot prototype with LiDAR-based mapping, localization, global planning, motion integration, semantic destinations, and a human-assisted elevator workflow.
  \item Open-sourced code, maps, configurations, and bilingual documentation.
\tightbulletsend

% ==================== Additional Projects ====================
\section{Additional Projects}

\compactentry
  {Exploratory RGB/NIR/Polarization Depth Estimation}
  {Aug. 2025--Oct. 2025}
  {Explored multimodal fusion, custom losses, knowledge distillation, calibration, and sparse LiDAR supervision.}

\entrygap

\compactentry
  {3D Dynamic Heart Reconstruction from 2D Medical Imaging}
  {Feb. 2025--May 2025}
  {Course Project Lead; led preprocessing, segmentation, reconstruction, and myocardial-motion visualization.}

\entrygap

\compactentry
  {Intelligent Risk Butler: Exchange-Rate Risk Inspection System}
  {Jan. 2025--Jun. 2025}
  {Core Member, 20th Citi Cup Algorithm Group; built an LLM module and exchange-rate risk signal pipeline.}

\entrygap

\compactentry
  {Quantum Implementation of Finite Element Analysis}
  {Jan. 2025--Feb. 2025}
  {Exchange Project, North Carolina State University; implemented block-encoding/QSP-based solvers and explored BPX preconditioning.}

\entrygap

\compactentry
  {Reservoir-Computing-Transistor-Based Fall Detection}
  {Nov. 2023--Dec. 2024}
  {National Undergraduate Innovation Program; contributed to a hardware-software system integrating analog computing, Arduino, Bluetooth, and vision-based motion analysis.}

% ==================== Internship ====================
\section{Internship Experience}

\researchentry
  {Data Analysis Intern}
  {Jul. 2024--Sep. 2024}
  {College of Water Conservancy and Architectural Engineering, Shihezi University}
\tightbullets
  \item Processed multivariate environmental sensor data and evaluated machine-learning models for anomaly detection in irrigation leakage and drought stress.
\tightbulletsend

% ==================== Skills ====================
\section{Technical Skills}
{\fontsize{11}{14.3}\selectfont
\renewcommand{\arraystretch}{1.06}
\hyphenpenalty=10000
\exhyphenpenalty=10000
\begin{tabularx}{\textwidth}{@{}>{\bfseries}p{1.40in}>{\raggedright\arraybackslash}X@{}}
Programming & Python, C, MATLAB, Arduino \\
Imaging \& Vision & RGB--NIR fusion, depth/pose estimation, feature matching, illumination estimation, geometric evaluation \\
Robotics & ROS 2, Nav2, visual odometry/SLAM, modular perception, Unitree Go2 \\
Sensors \& Data & RGB, active NIR, hyperspectral, polarization, LiDAR \\
Tools & PyTorch, OpenCV, NumPy, Linux, Git, LaTeX
\end{tabularx}\par}

\end{document}
